\section{Literature Positioning}
\label{sec:literature}

The formal structure employed in this preprint follows a well-established and
widely used paradigm in the mathematical treatment of observation and
identification.
A system is characterized by a state space, observation is represented as a map
from states to an observation space, and identifiability corresponds to
injectivity of that map on a domain of interest.
This abstract structure underlies classical work in control theory and system
identification and serves as the baseline against which the present results are
positioned.

\subsection{Observation as a Map and Loss of Information}

Modeling observation as a map
\(\pi:\Omega\to\Sigma\) and analyzing the induced equivalence relation on the
state space is classical.
Observation partitions \(\Omega\) into fibers (preimages) of \(\pi\), and any two
states lying in the same fiber are observationally indistinguishable under the
given observation scheme.
This perspective is already implicit in early state-space formulations of
control systems and observability
\cite{kalman1960}.

In control-theoretic and systems-theoretic language, this corresponds to the
standard distinction between \emph{observability} and
\emph{identifiability}.
Observability concerns whether outputs are defined and informative in principle,
whereas identifiability is the stronger property that states or parameters are
uniquely determined by outputs under specified conditions.
The fact that observability alone does not guarantee identifiability is
well-known and widely recognized in the literature.

\subsection{Identifiability in System Identification}

The system identification literature traditionally distinguishes between two
notions of identifiability.
\emph{Structural identifiability} refers to uniqueness under idealized,
noise-free observations within a given model class, while
\emph{practical identifiability} concerns uniqueness under finite data, noise,
and experimental limitations.
This distinction is classical and extensively documented
\cite{ljung1999,walter1997}.

The results of this preprint operate entirely at the structural level.
They formalize the fact that, within the Ontology of Continua and the frozen
executable specification \texttt{ETS.K\_EA.v1.1}, no general structural
identifiability guarantee follows for arbitrary finite KPI readout maps.
This is a statement about what is \emph{not implied} by the modeling framework
itself, independent of estimation algorithms, noise models, or data volume.
In this respect, the results are fully aligned with standard identifiability
theory rather than in tension with it.

\subsection{OC/ETS-Specific Contribution}

The only nonstandard element in the present analysis is the enterprise modeling
substrate.
OC and \texttt{ETS.K\_EA.v1.1} provide a formally specified enterprise continuum
with a state space \(\Omega(K_{EA})\), explicit structural components (axes,
flows, thresholds, cycles), and a factorized continuum measure.
In addition, ETS constrains admissible observation through explicit
observability conditions and invariants.

Within this substrate, KPI observation becomes a concrete instantiation of the
general observation-map paradigm:
\[
\pi_K:\Omega_{\mathrm{obs}}\to\Sigma_{\mathrm{obs}}\subseteq\mathbb{R}^m.
\]
Once this structure is fixed, the non-identifiability results follow directly,
without importing auxiliary assumptions such as smoothness, dimensionality
arguments, or genericity conditions.
The contribution of this preprint is therefore not a new impossibility theorem,
but the explicit localization of standard identifiability limits within the
OC/ETS enterprise continuum framework.

\subsection{Relation to Partial Observability and Representation Limits}

In many applied domains, including control, econometrics, and organizational
measurement, only partial projections of the underlying system are observable.
KPI dashboards are an archetypal example: they compress complex enterprise states
into a finite collection of scalar outputs.

From this perspective, indistinguishability of distinct underlying states is not
an anomaly but an expected consequence of projection.
Identifiability typically requires additional structure, such as richer sensing,
persistent excitation, experimental intervention, restrictive model classes,
or explicit injectivity constraints
\cite{slotine1991,walter1982}.

OC/ETS makes this separation explicit by distinguishing:
\begin{itemize}[leftmargin=2em]
  \item admissibility of observation via the observable domain
        \(\Omega_{\mathrm{obs}}\),
  \item the readout mechanism \(\pi_K\),
  \item and the uniqueness question captured by injectivity and
        identifiability.
\end{itemize}

The present note should therefore be read as a formal reminder that
\emph{projection-induced ambiguity is intrinsic} and does not disappear through
aggregation, smoothing, or deterministic post-processing alone.

\subsection{Positioning of the Present Results}

We draw on classical literature in observability, system identification, and
identifiability to anchor the formal notions used here
\cite{kalman1960,ljung1999,walter1997}.
No attempt is made to provide a comprehensive survey.
The purpose of this section is positioning rather than comparison: the logical
structure of the results is standard, while their application to KPI-based
enterprise observation within the OC/ETS continuum framework is specific and
diagnostic in nature.
